\documentclass[11pt]{article}
\usepackage[margin=1in]{geometry}
\usepackage[T1]{fontenc}
\usepackage{lmodern}
\usepackage{microtype}
\usepackage{xcolor}
\usepackage{hyperref}
\usepackage{booktabs}
\usepackage{enumitem}
\usepackage{listings}
\usepackage{tcolorbox}
\usepackage{parskip}

\hypersetup{
  colorlinks=true,
  linkcolor=blue!60!black,
  urlcolor=blue!60!black
}

\definecolor{codebg}{RGB}{245,245,245}
\definecolor{codecomment}{RGB}{90,100,90}
\definecolor{codekeyword}{RGB}{0,70,140}
\definecolor{codestring}{RGB}{150,60,20}

\lstdefinestyle{pythonstyle}{
  language=Python,
  backgroundcolor=\color{codebg},
  basicstyle=\ttfamily\small,
  keywordstyle=\color{codekeyword}\bfseries,
  commentstyle=\color{codecomment}\itshape,
  stringstyle=\color{codestring},
  showstringspaces=false,
  breaklines=true,
  frame=single,
  rulecolor=\color{black!20},
  columns=fullflexible,
  keepspaces=true
}
\lstdefinestyle{shellstyle}{
  language=bash,
  backgroundcolor=\color{codebg},
  basicstyle=\ttfamily\small,
  showstringspaces=false,
  breaklines=true,
  frame=single,
  rulecolor=\color{black!20},
  columns=fullflexible,
  keepspaces=true
}
\lstset{style=pythonstyle}

\newtcolorbox{keyidea}{
  colback=blue!5, colframe=blue!50!black,
  title=Key Idea, fonttitle=\bfseries
}
\newtcolorbox{tipbox}{
  colback=green!5, colframe=green!45!black,
  title=Tip, fonttitle=\bfseries
}
\newtcolorbox{warningbox}{
  colback=orange!7, colframe=orange!60!black,
  title=Common Pitfall, fonttitle=\bfseries
}

\title{
  \textbf{Codebase Change Summary --- Week 4}\\[4pt]
  \large Test-Driven Development, Mocking, and Dependency Isolation\\[2pt]
  \normalsize SE 413/513 --- Software Engineering
}
\author{Demo Testing Suite}
\date{}

\begin{document}
\maketitle

\begin{center}
\textit{What changed in \texttt{demo-testing-suite/} since Week 3, and why}
\end{center}

\vspace{0.3cm}

% ============================================================
\section{Purpose of This Document}
% ============================================================
The \texttt{demo-testing-suite/} codebase is a single, shared project that
grows one lecture at a time rather than being rebuilt from scratch each
week. This document summarizes exactly what was added or changed for
Week 4 and, more importantly, explains \emph{why} each change exists ---
what testing concept from the Week 4 reading it is meant to demonstrate.
Read this alongside the reading and \texttt{demo-testing-suite/README.md}.

\begin{keyidea}
Nothing from Weeks 1--3 was removed or rewritten. Everything below is
\emph{additive} --- the running examples you already saw
(\texttt{calculate\_late\_fee}, the shopping cart, discounts,
\texttt{determine\_shipping\_cost}) are all still there, unchanged, and
still behave the same way, including their two known, intentional
\texttt{xfail}/\texttt{FAIL} results.
\end{keyidea}

% ============================================================
\section{What Changed At a Glance}
% ============================================================
\begin{center}
\begin{tabular}{@{}p{5.8cm}p{1.6cm}p{7.0cm}@{}}
\toprule
\textbf{File} & \textbf{Status} & \textbf{Purpose} \\
\midrule
\texttt{app/notifications/password\_reset.py} & New & A direct
transcription of the reading's \texttt{PasswordResetService} worked
example: a dependency-injection seam (\texttt{email\_client}) replacing
a hard-coded \texttt{smtplib} dependency. \\
\texttt{tests/notifications/test\_password\_reset\_service.py} & New &
The reading's full dummy/stub/fake/spy/mock taxonomy, plus the
``mocks that lie'' pitfall, as seven runnable \texttt{pytest} tests. \\
\texttt{app/accounts/registration\_service.py} & New &
\texttt{UserRegistrationService}: the first caller of
\texttt{UserRepository} that needs isolating --- fulfills that file's
own Week 2 docstring forward-reference to a future mocking lecture. \\
\texttt{tests/accounts/test\_registration\_service.py} & New & Mocks
\texttt{UserRepository} directly (via \texttt{create\_autospec}) instead
of a real \texttt{sqlite3} connection, plus a second ``mocks that lie''
demonstration. \\
\texttt{docs/DEMO\_GUIDE\_week04\_tdd\_test\_doubles.md} & New & A
section-by-section lecture companion (primarily for in-class use). \\
\texttt{README.md} & Modified & Updated project layout description and
the expected \texttt{pytest -v} output, now that Week 4's files exist. \\
\bottomrule
\end{tabular}
\end{center}

Note that \texttt{tests/practice/tdd\_password\_demo.py} (the
Red-Green-Refactor live-coding script for the reading's
\texttt{is\_strong\_password()} example) was already added to the
repository ahead of this reading, in anticipation of this lecture. It is
unchanged by this update and is described here for completeness only.

% ============================================================
\section{New App Code: \texttt{PasswordResetService}}
% ============================================================
\textbf{Location:} \texttt{app/notifications/password\_reset.py}

This file is the centerpiece worked example from the reading's
``Dependency Isolation and Seams'' section, transcribed so it can be run
against real code instead of a slide. It defines three things:

\begin{itemize}
\item \textbf{\texttt{EmailClient}} --- an interface with no real
behavior. Its only job is to give \texttt{unittest.mock.create\_autospec}
something concrete to constrain a mock to (see the Common Pitfalls
tests, below).
\item \textbf{\texttt{SmtpEmailClient}} --- the real, production
implementation, wrapping \texttt{smtplib}. It is intentionally
\emph{not} exercised by any test in this suite --- testing it for real
would require a real mail server, which is exactly the problem
dependency injection solves.
\item \textbf{\texttt{PasswordResetService}} --- the class actually
under test. Its constructor accepts \texttt{email\_client} and
\texttt{logger} from outside instead of constructing them internally.
\end{itemize}

\begin{keyidea}
The reading's ``before'' version --- a hard-coded
\texttt{smtplib.SMTP(...)} call built directly inside
\texttt{send\_reset\_email()} --- is described in this file's docstring
but deliberately \emph{not} reproduced as runnable code. Only the
``after'' version, with the seam already in place, actually exists here.
This mirrors how Week 3's \texttt{shipping.py} carried only the
specification, not a second broken copy of the function.
\end{keyidea}

% ============================================================
\section{New Tests: The Test Double Taxonomy}
% ============================================================
\textbf{Location:} \texttt{tests/notifications/test\_password\_reset\_service.py}

This file encodes the reading's five test doubles as seven runnable
tests, in the same order the reading introduces them:

\begin{itemize}
\item \textbf{Dummy} --- \texttt{DummyLogger}, poisoned so any
accidental use raises immediately. \texttt{send\_reset\_email()} never
calls \texttt{logger} at all, so a correct test never triggers it.
\item \textbf{Stub} --- \texttt{StubEmailClient}, a canned answer with
no assertion made about how it was called.
\item \textbf{Fake} --- \texttt{FakeEmailClient}, verified via
\textbf{state verification} (inspecting \texttt{sent\_emails}
afterward).
\item \textbf{Spy} --- \texttt{SpyEmailClient}, which records calls for
the test to inspect.
\item \textbf{Mock} --- a bare \texttt{MagicMock()}, verified via
\textbf{behavior verification}
(\texttt{assert\_called\_once\_with(...)}).
\end{itemize}

The file closes with two tests demonstrating the reading's ``Mocks that
lie'' pitfall directly: a bare \texttt{MagicMock()} silently accepts a
typo'd method name (\texttt{sedn} instead of \texttt{send}), while
\texttt{create\_autospec(EmailClient, instance=True)} raises an
\texttt{AttributeError} on the identical typo.

\begin{tipbox}
Run the suite yourself:
\begin{lstlisting}[style=shellstyle]
cd demo-testing-suite
pytest tests/notifications/test_password_reset_service.py -v
\end{lstlisting}
All seven tests pass. Unlike Weeks 1 and 3, there is no deliberate bug
to find in this file --- the teaching payload is in reading the five
doubles side by side, not in a failure.
\end{tipbox}

% ============================================================
\section{New App Code and Tests: Mocking \texttt{UserRepository}}
% ============================================================
\textbf{Locations:} \texttt{app/accounts/registration\_service.py},
\texttt{tests/accounts/test\_registration\_service.py}

\texttt{app/accounts/user\_repository.py}'s docstring has said, since
Week 2: ``\emph{A mocking lecture can mock UserRepository instead of
sqlite3 itself.}'' Until now, nothing in the codebase actually called
\texttt{UserRepository}, so that sentence had nothing to point at.
\texttt{UserRegistrationService} is that first caller: it accepts a
\texttt{user\_repository} and an \texttt{email\_client} via constructor
injection, registers a new user through the repository, and welcomes
them through the email client.

Its test file never touches \texttt{sqlite3}. Instead,
\texttt{create\_autospec(UserRepository, instance=True)} replaces the
whole dependency, and the tests check both outcomes (a returned user id,
a rejected duplicate) and interactions (the welcome email's exact
arguments, or that no email is sent on rejection). A third test repeats
the ``mocks that lie'' point from the notifications suite, this time
against \texttt{UserRepository}'s own interface.

\begin{tipbox}
Contrast this with \texttt{tests/accounts/test\_user\_repository.py},
which still exercises a real \texttt{:memory:} SQLite connection ---
that file tests \texttt{UserRepository} itself, while this new file
tests a \emph{caller} of \texttt{UserRepository} in isolation from it.
\begin{lstlisting}[style=shellstyle]
pytest tests/accounts/test_registration_service.py -v
\end{lstlisting}
\end{tipbox}

% ============================================================
\section{New Documentation}
% ============================================================
\textbf{Location:} \texttt{docs/DEMO\_GUIDE\_week04\_tdd\_test\_doubles.md}

This is a lecture-companion document, primarily written to guide what
gets shown and run during class, section by section, following the
Week 4 reading's own structure --- including how to live-code
\texttt{tests/practice/tdd\_password\_demo.py}'s Red-Green-Refactor
cycles. You are welcome to read it, but it is not required reading ---
the reading itself is the primary material.

% ============================================================
\section{Updated: \texttt{README.md}}
% ============================================================
The top-level \texttt{demo-testing-suite/README.md} was updated to
reflect the four new files above and to describe the suite's current
expected output. If you run the full suite from the project root:

\begin{lstlisting}[style=shellstyle]
cd demo-testing-suite
pip install -r requirements.txt
pytest -v
\end{lstlisting}

you should now see \textbf{43 passed, 2 xfailed, 1 failed} --- ten more
passing tests than Week 3's baseline, and the same two \texttt{xfail}
results and one \texttt{FAIL} as before, all unrelated to this week's
additions.

% ============================================================
\section{How This Connects to the Reading}
% ============================================================
\begin{keyidea}
The reading argues that TDD's real value is design pressure: writing a
test first surfaces hidden dependencies before they calcify. Dependency
injection turns a hidden dependency into a seam, and the dummy/stub/
fake/spy/mock vocabulary gives a precise way to choose what belongs on
the other side of that seam during a test.
\texttt{app/notifications/password\_reset.py} and
\texttt{app/accounts/registration\_service.py} exist so you can see both
halves of that argument against real, running code: a seam introduced
where a hard-coded dependency used to be, and all five kinds of double
filling it in five different tests.
\end{keyidea}

% ============================================================
\section{Suggested Next Steps}
% ============================================================
\begin{enumerate}
\item Read \texttt{app/notifications/password\_reset.py}'s docstring and
identify, in your own words, why the reading's ``before'' version would
be hard to unit test --- before looking at how the ``after'' version
solves it.
\item Run \texttt{pytest tests/notifications/test\_password\_reset\_service.py -v}
and, for each of the five doubles, say out loud which school (classical
or mockist) it belongs to and why.
\item Run \texttt{pytest tests/accounts/test\_registration\_service.py -v}
and explain why this file needs no real database at all, in contrast to
\texttt{tests/accounts/test\_user\_repository.py}.
\item Try Exercise 2 from the reading (\texttt{SessionManager} /
\texttt{time.time()}) on paper: identify the hard-coded dependency and
sketch the constructor-injected version, the same way this week's app
code did for \texttt{PasswordResetService} and
\texttt{UserRegistrationService}.
\end{enumerate}

\end{document}
